%% kapa: exam=1 round=36 year=2025 period=1 form=A subject=civil_law subject_ko=민법 seq=32 number=32
%% source: Q-Net 감정평가사 기출문제 (cst00309) artlSeq=5231525 file=raw/1008/5231525/2025년 제36회 감정평가사 1차시험 1교시 문제지 원본.pdf
%% conversion: kordoc 4.24.51 -> tools/gen1.py (auto); review=proofread-sonnet; flags=none
\begin{question}{32}
관습상의 법정지상권이 성립할 수 있는 경우는? (다툼이 있으면 판례에 따름)
\begin{choices}
  \item 자신의 토지에 미등기건물을 건축한 자가 건물의 철거특약 없이 토지를 매도한 경우
  \item 토지소유자가 토지매도 후 소유권이전등기 전에 매수인의 승낙 없이 그 토지에 건물을 신축한 경우
  \item 건물에 대한 경매개시결정 후 매수인의 매각대금완납 전에 그 건물과 대지가 동일인에게 속하게 된 경우
  \item 토지공유자 1인이 지분의 과반수의 동의를 얻어 공유토지 위에 건물을 신축하여 소유하다가 그 건물을 타인에게 매도한 경우
  \item 채권담보를 위한 가등기가 설정된 나대지 위에 건물이 신축된 후 그 가등기에 기한 본등기가 행해진 경우
\end{choices}
\end{question}
\answer{1}
